\begin{titlepage}

\begin{figure}[H]
  \begin{center}
    \includegraphics[width=3cm]{auth.pdf}
    \label{fig:cover_auth_logo}
  \end{center}
\end{figure}

\centering
\Large Αριστοτέλειο Πανεπιστήμιο Θεσσαλονίκης\\
\Large Πολυτεχνική Σχολή\\
\large Τμήμα Ηλεκτρολόγων Μηχανικών και Μηχανικών Υπολογιστών\\
\large Τομέας Ηλεκτρονικής \& Υπολογιστών

\vspace{\fill}

\LARGE Ψηφιακά Δίδυμα στην Εκπαίδευση: Ανάπτυξη Συστήματος \en{IoT} για τη Βελτιστοποίηση της Ενεργειακής Απόδοσης και του Μαθησιακού Περιβάλλοντος

\vspace{\fill}

\Large Διπλωματική Εργασία\\
\Large του\\
\Large Αλέξανδρου Παπαδόπουλου

\vspace{\fill}
\raggedright

\begin{tabular}{ll}
\textbf{Επιβλέπων:} & Αλκιβιάδης Χατζόπουλος\\
\end{tabular}

\centering
\vspace{\fill}
\today

\end{titlepage}

\chapter*{Περίληψη}
\addcontentsline{toc}{chapter}{Περίληψη}
Η παρούσα διπλωματική εργασία πραγματεύεται τον σχεδιασμό, την υλοποίηση και την αξιολόγηση μιας ολοκληρωμένης
πλατφόρμας Ψηφιακού Διδύμου (\en{Digital Twin}), ειδικά προσαρμοσμένης για εκπαιδευτικά περιβάλλοντα. Σε αντίθεση
με τις παραδοσιακές, αντιδραστικές (\en{reactive}) μεθόδους διαχείρισης κτιρίων, η προτεινόμενη λύση εγκαθιδρύει μια
αμφίδρομη, πραγματικού χρόνου επικοινωνία μεταξύ του φυσικού σχολικού χώρου και του ψηφιακού του αντιγράφου. Η
ανάπτυξη του συστήματος ακολούθησε μια αυστηρή διεπιστημονική προσέγγιση. Σε επίπεδο υλικού (\en{Hardware}),
εγκαταστάθηκε δίκτυο \en{IoT} χαμηλού κόστους (μικροελεγκτές \en{ESP32}, έξυπνοι μετρητές \en{Shelly} και συσκευές
\en{Sonoff}). Σε επίπεδο λογισμικού (\en{Software}), αναπτύχθηκε μια ισχυρή αρχιτεκτονική (\en{MQTT}, \en{InfluxDB},
\en{openHAB}), η οποία τροφοδοτεί μια προσαρμοσμένη διεπαφή χρήστη (\en{React 3D Dashboard}) για την οπτικοποίηση
των δεδομένων. Η έρευνα δομείται πάνω σε τρεις κεντρικούς πυλώνες: (α) τη βελτιστοποίηση του Μαθησιακού
Περιβάλλοντος (\en{Environmental Quality}), μέσω της παρακολούθησης της ακουστικής άνεσης για την ενίσχυση της
συγκέντρωσης των μαθητών, (β) την Ενεργειακή Αποδοτικότητα (\en{Energy Efficiency}), με συνεχή καταγραφή ισχύος,
και (γ) τη Λειτουργική Διαχείριση (\en{Facility Management}) μέσω κεντρικού ελέγχου. Από την πιλοτική εφαρμογή στα
«Εκπαιδευτήρια Πλάτων», επιτεύχθηκε και επαληθεύτηκε μέση μείωση της ενεργειακής κατανάλωσης κατά σχεδόν 20\%,
επιβεβαιώνοντας μέσω τεχνοοικονομικής ανάλυσης την υψηλή προστιθέμενη αξία της μετάβασης σε ένα σύγχρονο, Έξυπνο
Σχολικό Συγκρότημα (\en{Smart Campus}).

\textbf{Λέξεις-Κλειδιά:} Ψηφιακό Δίδυμο, \en{Smart Campus}, Διαδίκτυο των Πραγμάτων (\en{IoT}), Μαθησιακό
Περιβάλλον, Ενεργειακή Αποδοτικότητα, Ποιότητα Εσωτερικού Περιβάλλοντος (\en{IEQ}).

\chapter*{\en{Abstract}}
\addcontentsline{toc}{chapter}{\en{Abstract}}
\selectlanguage{english}
This diploma thesis presents the design, implementation, and evaluation of a comprehensive Digital Twin platform
tailored for educational environments. Shifting away from traditional, reactive facility management methods, the
proposed solution establishes a bidirectional, real-time communication channel between the physical school environment
and its digital counterpart. The system's development followed a rigorous interdisciplinary approach. On the hardware
level, a network of low-cost IoT devices (ESP32 microcontrollers, Shelly smart meters, and Sonoff actuators) was
deployed. On the software level, a robust architecture (MQTT, InfluxDB, openHAB) was developed to feed a custom-built
React 3D Dashboard for immersive data visualization. The core of this research is built upon three central pillars:
(a) Environmental Quality, optimizing the learning environment through strict acoustic monitoring to improve student
focus; (b) Energy Efficiency, featuring real-time power tracking; and (c) Facility Management, enabling centralized
control of infrastructure. Following the pilot implementation at "Platon Schools," a verified reduction in energy
consumption of approximately 20\% was recorded. The accompanying techno-economic analysis confirmed the high added
value, financial viability, and structural necessity of transitioning towards a modern Smart Campus ecosystem.

\textbf{Keywords:} Digital Twin, Smart Campus, Internet of Things (\en{IoT}), Learning Environment, Energy
Efficiency, Indoor Environmental Quality (\en{IEQ}).
\selectlanguage{greek}

\chapter*{Ευχαριστίες}
\addcontentsline{toc}{chapter}{Ευχαριστίες}
Η παρούσα διπλωματική εργασία αποτελεί το επιστέγασμα μιας μακράς, απαιτητικής, αλλά και εξαιρετικά δημιουργικής
προσπάθειας. Η ολοκλήρωσή της δεν θα ήταν εφικτή χωρίς την πολύτιμη συμβολή, καθοδήγηση και υποστήριξη ανθρώπων
που στάθηκαν δίπλα μου και πίστεψαν στο όραμα αυτής της πλατφόρμας.

Αρχικά, θα ήθελα να εκφράσω τις θερμότερες και ειλικρινείς ευχαριστίες μου στον επιβλέποντα καθηγητή μου,
κ. Άλκη Χατζόπουλο, για την εμπιστοσύνη που μου έδειξε, την πολύτιμη ακαδημαϊκή του καθοδήγηση και τη συνεχή
υποστήριξη καθ' όλη τη διάρκεια εκπόνησης της εργασίας.

Ένα ξεχωριστό και μεγάλο ευχαριστώ οφείλω στον κ. Χρήστο Κάλλα. Η συμβολή του υπήρξε καθοριστική, καθώς
αποτέλεσε την αφετηρία και την τεχνική βάση ολόκληρης της έρευνας. Η βοήθειά του στην εγκατάσταση του υλικού
εξοπλισμού (\en{hardware}) και των αισθητήρων \en{IoT}, καθώς και η πολύτιμη καθοδήγησή του στο στήσιμο και την
παραμετροποίηση της πλατφόρμας \en{openHAB}, αποτέλεσαν το θεμέλιο πάνω στο οποίο αναπτύχθηκε το παρόν Ψηφιακό
Δίδυμο.

Τέλος, θα ήθελα να ευχαριστήσω θερμά τη διοίκηση, τους εκπαιδευτικούς και το προσωπικό των «Εκπαιδευτηρίων
Πλάτων» στην Κατερίνη. Η προθυμία τους να αποδεχθούν το σχολείο τους ως πεδίο μελέτης περίπτωσης
(\en{case study}) και να φιλοξενήσουν την πιλοτική εφαρμογή του συστήματος, μου έδωσε τη μοναδική ευκαιρία να
εφαρμόσω τη θεωρία στην πράξη σε ένα πραγματικό, ζωντανό εκπαιδευτικό περιβάλλον.

\clearpage
